%-----------------------------------------------------------
%   Capítulo - Projeto
%-----------------------------------------------------------
\chapter{Projecto}
\label{cap:projecto}

Este capítulo transforma a arquitetura de alto nível do
Capítulo~\ref{cap:arquitectura} numa solução concreta. Fecham-se as
decisões de projeto DP01--DP17 que dependem da escolha de componentes e da
organização do software, comparam-se alternativas com preços e
apresentam-se duas abordagens de implementação com o respetivo orçamento.
Os preços indicados são preços unitários sem IVA, consultados em
distribuidores europeus em novembro de 2026.

\section{Arquitectura do Sistema}
\label{sec:proj-arquitetura}

A Figura~\ref{fig:blocos-projeto} apresenta o diagrama de blocos da
solução. O microcontrolador concentra o processamento, o encaminhamento e
a referência temporal. As interfaces que ligam ao PC ficam à esquerda e
as que ligam a dispositivos e cargas à direita.

\begin{figure}[htbp]
\centering
\Diagrama{blocos_projeto}
\caption[Diagrama de blocos da solução.]{Diagrama de blocos da solução, com os componentes selecionados. As interfaces isoladas têm alimentação isolada integrada.}
\label{fig:blocos-projeto}
\end{figure}

Os dois barramentos CAN FD de RF04 ficam com papéis distintos (DP01): um
dedicado à ligação ao PC, como alternativa de RF09, e outro a nós
externos. Quando o PC usa Ethernet ou RS232, o primeiro pode também
servir nós externos, por configuração. A Ethernet liga apenas ao PC.

Quanto à proteção das comunicações (DP10), isolam-se galvanicamente os
barramentos que saem da unidade para nós remotos (CAN FD e RS485) e as
saídas PWM, onde o isolamento é obrigatório (RF15). RS232 e I2C são
ligações locais e ficam sem isolamento, com proteção contra ESD.

\section{Opções de Projecto}
\label{sec:proj-opcoes}

\subsection{Unidade de processamento}

A Tabela~\ref{tab:proj-mcu} compara os microcontroladores avaliados.
Os quatro cumprem os requisitos de comunicação com, pelo menos, dois
controladores CAN FD e um MAC Ethernet.

\begin{table}[htbp]
\centering\small\setstretch{1.05}
\renewcommand{\arraystretch}{1.15}
\caption{Comparação dos microcontroladores candidatos.}
\label{tab:proj-mcu}
\begin{tabularx}{\textwidth}{@{}>{\raggedright\arraybackslash}p{0.2\textwidth}>{\raggedright\arraybackslash}X>{\raggedright\arraybackslash}X>{\raggedright\arraybackslash}X>{\raggedright\arraybackslash}X@{}}
\toprule
\rowcolor{LasetTint}
Critério & STM32H563ZI & STM32H743ZI & SAM E70Q21 & i.MX RT1176 \\
\midrule
Núcleo & M33, 250 MHz & M7, 480 MHz & M7, 300 MHz & M7 1 GHz + M4 \\
CAN FD / Ethernet & 2 / 10/100 & 2 / 10/100 & 2 / 10/100 & 3 / 10/100 + 2 Gbit \\
Flash interna & 2 MB, 2 bancos & 2 MB, 2 bancos & 2 MB & externa \\
Encapsulamento & LQFP144 & LQFP144 & LQFP144 & BGA289 \\
Preço do chip & 7,68 € & 12,83 € & 16,19 € & 16,26 € \\
Placa de avaliação & 25,67 € & 34,68 € & 156,87 € & 164,83 € \\
\bottomrule
\end{tabularx}
\end{table}

O i.MX RT1176 foi excluído por existir apenas em BGA289, sem flash
interna, e por o número de peça estar marcado como não recomendado para
novos projetos. O SAM E70 cumpre os requisitos, mas a placa de avaliação
custa cerca de seis vezes mais do que as da ST. Entre os dois STM32, o
STM32H563 é mais barato, mas a pilha TCP/IP oficial da ST para esta
família é a NetX Duo sobre ThreadX e o LwIP com FreeRTOS exige porte
manual. Escolheu-se o \textbf{STM32H743ZIT6} (DP03): o FreeRTOS, candidato
indicado pelo docente, e o LwIP estão integrados no STM32CubeMX para esta
família, com exemplos oficiais em placas Nucleo~\cite{st_lwip_h7}. Os seus
controladores FDCAN atingem 8~Mbit/s na fase de dados, com a fase de
arbitragem limitada a 1~Mbit/s~\cite{st_fdcan_8m}, o que é coerente com a
nota de RF04 sobre as duas fases.

\subsection{Interfaces e alimentação}

A Tabela~\ref{tab:proj-componentes} resume os componentes escolhidos para
cada bloco, com a alternativa considerada.

\begin{table}[htbp]
\centering\small\setstretch{1.05}
\renewcommand{\arraystretch}{1.15}
\caption{Componentes selecionados por bloco.}
\label{tab:proj-componentes}
\begin{tabularx}{\textwidth}{@{}>{\raggedright\arraybackslash}p{0.17\textwidth}>{\raggedright\arraybackslash}p{0.19\textwidth}>{\raggedright\arraybackslash}X>{\raggedright\arraybackslash}p{0.1\textwidth}@{}}
\toprule
\rowcolor{LasetTint}
Bloco & Componente & Justificação & Preço \\
\midrule
CAN FD ($\times2$) & ADM3057E & Isolado 3 kV com alimentação isolada integrada; até 12 Mbit/s; proteção de $\pm40$ V no barramento~\cite{adi_adm3057e}. & 13,67 € \\
RS485 & ADM2582E & Isolado 2,5 kV, 16 Mbit/s, até 256 nós, funciona a 3,3 V~\cite{adi_adm2582e}. & 13,97 € \\
RS232 & MAX3227E & Taxa de 1 Mbit/s garantida, de 3,0 a 5,5 V~\cite{maxim_max3227e}. & 5,96 € \\
PWM ($\times3$) & ADuM6401 & Três canais isolados de 5 kV mais um de retorno; alimentação isolada de 5 V integrada~\cite{adi_adum6401}. Alimentado a 5 V, recebe os PWM do MCU através de três 74AHCT1G125, que aceitam a lógica de 3,3 V. & 14,91 € \\
Ethernet & LAN8742A & PHY RMII usado nas placas Nucleo, compatível com o controlador do CubeMX. & 1,30 € \\
I2C ($\times2$) & PCA9306 & Tradução de nível entre 3,3 V e sensores de 1,8 a 5,5 V. & 0,65 € \\
Fonte de 5 V & LM5177 & Buck-boost de 3,5 a 60 V; mantém 5 V com 5 V à entrada~\cite{ti_lm5177}. & 6,45 € \\
Fonte de 3,3 V & TPS62162 & Buck síncrono de 1 A a partir dos 5 V. & 1,35 € \\
Polaridade & STPS5H100 & Schottky de 100 V e 5 A em série com a entrada. Substitui o díodo ideal LM74700-Q1, que só aguenta 65 V, abaixo da tensão de limitação do TVS de entrada. & a confirmar \\
\bottomrule
\end{tabularx}
\end{table}

A gama de entrada de RNF08 determina a topologia da fonte (DP13). Os
transceivers CAN FD e o isolador dos PWM precisam de 5 V, e um conversor
abaixador não consegue gerar 5 V quando a entrada também é de 5 V. Por
isso usa-se um conversor buck-boost, seguido de um abaixador para 3,3 V.

\subsection{Abordagens de implementação}

Comparam-se duas abordagens (DP05):
\begin{description}
    \item[A. Placa de avaliação e placa de interfaces:] \hfill\\*
    NUCLEO-H753ZI com uma placa de duas camadas que contém apenas os
    transceivers e os isoladores. É mais rápida de obter, mas não inclui
    a fonte de 5 a 60 V, pelo que não demonstra RNF08. Custo estimado dos
    componentes com preço confirmado: 163,18 €.
    \item[B. PCB própria:] \hfill\\*
    Placa de quatro camadas com o STM32H743, as interfaces e a fonte
    completa. Cumpre todos os requisitos elétricos. Custo estimado dos
    componentes com preço confirmado: 197,94 €.
\end{description}

Ambos os valores incluem a placa NUCLEO e um recetor GPS para os ensaios
de RF16, e excluem passivos, conectores e o fabrico da PCB, ainda por
fechar com o esquema elétrico. Propõe-se a abordagem B como solução
final, usando a NUCLEO para desenvolver o firmware enquanto a PCB é
fabricada.

\section{Descrição de \textit{Hardware}}
\label{sec:proj-hardware}

O esquema elétrico completo está no Anexo~\ref{anexo:EETS}. É gerado por
código a partir da pasta \texttt{hardware/gerador} do repositório, o que
permite rever cada alteração como o resto do relatório, e está dividido em
sete folhas: alimentação, microcontrolador, Ethernet, CAN FD, RS485 e RS232,
I2C e GPS, e PWM isolado.

A árvore de alimentação é a seguinte. A entrada de 5 a 60 V passa por um
fusível rearmável, um TVS e o díodo Schottky de proteção contra inversão.
O LM5177 gera o barramento de 5 V para os transceivers CAN FD e o isolador
dos PWM, com quatro MOSFET BSC160N10NS3 de 100 V, uma bobina de 18~$\mu$H
e comutação a 400 kHz. A malha de controlo foi compensada para uma largura
de banda de 4 kHz, abaixo de um terço do zero no semiplano direito do modo
elevador (14,5 kHz) com 2 A de carga. O TPS62162 gera os 3,3 V para o
microcontrolador, o PHY, o RS232, o RS485 e o lado lógico dos isoladores
CAN FD.

A atribuição de pinos do STM32H743 segue a da placa NUCLEO-H743ZI na
Ethernet RMII, para reaproveitar o projeto do CubeMX: FDCAN1 em PD0/PD1,
FDCAN2 em PB5/PB6, RS232 no USART2, RS485 no USART6 com controlo automático
do emissor (DE), GPS no UART7 e os três PWM nos canais 1 a 3 do TIM1. A
linha de falha do conector PWM entra pela entrada de paragem do TIM1
(BKIN), o que desliga as saídas por hardware sem depender do firmware. O
PPS do recetor GPS liga a PA0, canal de captura do TIM2, que é de 32 bits.

Ficam por fechar antes do fabrico: a margem de fase da fonte de 5 V,
medida no protótipo, o footprint do LM5177, feito a partir do encapsulamento
TSSOP-38 com o pad exposto, e a escolha dos conectores definitivos.

\section{Descrição do \textit{Software}}
\label{sec:proj-software}

O firmware executa sobre FreeRTOS com memória estática: filas e
\textit{pools} de blocos são criados no arranque. As rotinas de
interrupção e o DMA de cada interface apenas recebem dados e colocam-nos
numa fila de entrada. Uma tarefa de encaminhamento aplica a tabela de
encaminhamento configurada e coloca cada mensagem na fila de saída do
destino. Uma tarefa de transmissão por interface esvazia essa fila ao
ritmo do respetivo barramento, pelo que uma saída lenta não impede o
progresso das restantes (RF12).

Cada fila de saída tem duas classes, urgente e normal (RF14). As
mensagens urgentes são atendidas primeiro, mas, após um número
configurável de urgentes seguidas, passa uma normal, o que garante
progresso a ambas. O dimensionamento das filas (RF13, DP09) parte de
$B = (R_{\mathrm{in}} - R_{\mathrm{out}})\,T$, em que $T$ é a duração da
rajada, e é ajustado com os contadores de descarte nos ensaios.

A comunicação com o PC (DP04) usa um formato de mensagem comum, com
versão, tipo, interface, canal, endereço ou identificador, número de
sequência, marca temporal de 64 bits em microssegundos e comprimento.
Em Ethernet, as operações seguem por UDP, com confirmação na aplicação,
e a configuração por TCP. Em RS232, as mensagens são delimitadas com COBS
e CRC-16. Em CAN FD, as mensagens com mais de 64 bytes são segmentadas
segundo a ISO 15765-2.

A sincronização (RF16, DP14) usa um temporizador de 32 bits a 240 MHz
com captura no flanco do PPS, o que dá uma resolução de cerca de 4 ns. A
mensagem NMEA indica o segundo UTC e a frequência real do temporizador é
medida entre impulsos PPS sucessivos, compensando a deriva do cristal.
Sem PPS, a unidade mantém a última frequência medida e assinala que não
está sincronizada.

Os três canais PWM usam um temporizador avançado com contagem a 2 MHz,
o que dá uma resolução de 0,5~$\mu$s (DP12). No arranque, as saídas ficam em
nível baixo até existir uma configuração válida e, sem comandos do PC
durante um intervalo configurável, passam a um valor seguro.

A configuração (RF18, DP15) fica no segundo banco da flash interna, em
dois setores usados alternadamente. Cada registo tem um identificador,
a versão do formato, um número de sequência e um CRC32. No arranque
usa-se o registo válido mais recente e, se nenhum for válido, os valores
de fábrica. Como o código corre do primeiro banco, a escrita no segundo
não interrompe a execução.
